\chapter*{HISTORICAL NOTE}
\phantomsection
\addcontentsline{toc}{chapter}{HISTORICAL NOTE}

(N.B. --- The Roman numerals refer to the bibliography at the end of this note.)

The concepts of arithmetic mean and geometric mean of two positive quantities go back to antiquity, in particular the Pythagoreans made of them one of their favorite subjects. It is also probable that the inequality \(\sqrt{ab} \leqslant \frac{1}{2}(a + b)\) between these means was well-known to them; in any case it is proved by Euclid, in connection with the problem of maximizing the product of two numbers whose sum is given. Apparently one must wait until the rather late date of 1729 to find explicit mention, by MacLaurin, of the generalization of this problem to n numbers, and the corresponding inequality (and this, even though much more difficult extremum problems had been addressed long before).

Other analogous inequalities appeared in Analysis and Geometry from the end of the 18th century onward. Thus, in 1789, Lhuilier solved a geometric maximum problem associated with the inequality

\[
\sqrt {\left(\sum_ {k = 1} ^ {n} a _ {k}\right) ^ {2} + \left(\sum_ {k = 1} ^ {n} b _ {k}\right) ^ {2}} \leqslant \sum_ {k = 1} ^ {n} \sqrt {a _ {k} ^ {2} + b _ {k} ^ {2}};
\]

Cauchy, in 1821, proved the special case

\[
\left(\sum_ {k = 1} ^ {n} a _ {k} b _ {k}\right) ^ {2} \leqslant \left(\sum_ {k = 1} ^ {n} a _ {k} ^ {2}\right) \left(\sum_ {k = 1} ^ {n} b _ {k} ^ {2}\right)
\]

of Hölder's inequality; Buniakowsky in 1859 and H. A. Schwarz in 1885 generalized Cauchy's inequality to integrals.

Nevertheless, it was not until nearly the end of the 19th century that the inequalities of convexity became the object of systematic study. O. Hölder, in 1888 (I), had the idea of introducing convex functions of a single variable into the question: he thus obtained the inequality of the geometric mean starting with the concavity of \(\log x\) ; applying the same idea to the function \(x^{r}\) , he proved the inequality that bears his name, and which had already been found a year earlier by L. J. Rogers starting from the inequality of the geometric mean (cf.~Exer. 4). As for Minkowski's inequality, it was proved by the latter in 1896 (for finite sums), in connection with his memorable works on the `Geometry of numbers' ((II), pp.~115-117); however, while the idea of convexity (for functions of any number of variables) is one of the fundamental concepts of this work, it is rather surprising that Minkowski did not notice that his inequality may be obtained by Hölder's method applied to \((1 + x^{r})^{1 / r}\) (he limited himself to applying the classical methods for seeking an extremum by means of the infinitesimal calculus).

For a deeper study of the inequalities of convexity and of their applications, the reader may consult the book of Hardy, Littlewood and Pólya devoted to the question, which also contains a very complete bibliography (III).

\begin{enumerate}
\def\labelenumi{(\Roman{enumi})}
\item
  O. HÖLDER, Ueber einen Mittelwertsatz, Göttinger Nachrichten (1889), pp.~38-47.
\item
  H. MINKOWSKI, Geometrie der Zahlen, 2nd. edn., Leipzig--Berlin (Teubner), 1910.
\item
  G. H. HARDY, J. E. LITTLEWOOD, G. PÓLYA, Inequalities, Cambridge (University Press), 1934.
\end{enumerate}

